\documentclass[12pt]{article}

\usepackage[margin=1in]{geometry}
\usepackage{amsmath, amsthm}
\usepackage{amssymb}
\usepackage{graphicx}
\usepackage{times}
% \usepackage{subcaption}
\usepackage[colorlinks=true, linkcolor=black, citecolor=black, urlcolor=black]{hyperref}


\title{CSCI 5352: Homework 5}
\author{
    Zachary Caterer % $^{1,2,3}$ \\
    % \small $^1$Department of Computer Science, University of Colorado Boulder \\
    % \small $^2$Department of Chemical and Biological Engineering, University of Colorado Boulder \\
    % \small $^3$Interdisciplinary Quantitative Biology Program, University of Colorado Boulder
}
\date{\today}

\begin{document}

\maketitle

\section*{Problem 3}

I chose the paper titled ``Quantifying hierarchy and dynamics in US faculty hiring and retention'' by Wapman et al., published in Nature in October 2022. The paper presents a comprehensive analysis of faculty hiring and retention across all PhD-granting institutions in the United States over the decade from 2011 to 2020. I choose this paper due to its overlap with my project for this class. 

The central research question posed by the authors was: How do patterns of hiring and attrition among tenure-track faculty in the U.S. reflect underlying structures of prestige, inequality, and demographic change? More specifically, they aimed to understand how institutional prestige shapes the academic workforce and how factors such as gender, field, and international status influence faculty retention.

To address this question, the authors curated and analyzed an extensive dataset comprising over 295,000 faculty members across 10,612 departments. They constructed directed networks representing faculty flows between doctoral and hiring institutions, measured attrition and self-hiring rates, and examined disparities across gender and international lines. Their network-based approach enabled the quantification of prestige hierarchies and the identification of systemic inequalities within and across disciplines.

The authors did an excellent job of collecting a massive, well-structured dataset and applying rigorous network science techniques to reveal broad, empirically grounded patterns. They also clearly articulated the implications of their findings for academic policy and diversity initiatives. Their method of aggregating field-level data into discipline-wide and institution-wide insights was both innovative and impactful.

One limitation of the study lies in the binary classification of gender, which, although acknowledged by the authors, excludes non-binary identities. Additionally, while the analysis is statistically robust, it remains largely descriptive and stops short of exploring causal mechanisms behind the observed inequalities.

Possible extensions of this work include incorporating non-binary and more nuanced identity markers, analyzing how hiring and attrition patterns have changed post-2020 (especially in response to the COVID-19 pandemic), and investigating how these dynamics play out in relation to race, ethnicity, and socioeconomic background. Another valuable direction would be to model intervention strategies and simulate their long-term impacts on academic diversity and equity.


\end{document}